%=============================================================================
%  06_applications.tex -- Tested predictions
%
%  CONVENTION USED IN THIS FILE:
%    * Every MEASURED quantity is a macro from macros.tex. Where macros.tex
%      does not yet define one, \needsnum{suggestedMacroName} appears in its
%      place; the value is in paper/NUMBERS.md and must be promoted to a
%      macro before submission. No measured value is written as a literal.
%    * Experimental *design* parameters -- sweep points, chain depths, schema
%      widths, deadline settings, block sizes, sample sizes -- are written as
%      literals, because they are settings we chose rather than results we
%      measured.
%=============================================================================
\section{Tested Predictions}
\label{sec:applications}

The two cost laws of \cref{sec:cost-laws} and the batching mechanism of
\cref{sec:mechanism} are not descriptions of a benchmark; they are claims about
which \emph{workload shapes} favour which architecture. A cost law that is flat
in output length predicts an advantage wherever output length is large or
variable. A batching multiplier of ${\sim}1\times$ predicts that the advantage
survives only where the competing architecture cannot deploy its own multiplier
--- that is, where the workload is serial by construction, by regulation, or by
deployment shape.

This section turns those two statements into five falsifiable predictions and
measures each one. All five experiments run on Campaign~B hardware
(\cref{sec:repro}) unless stated otherwise: \modelDG{} served by vLLM at
\texttt{--max-num-seqs 1}, against \modelARthree{} served by the same engine.
We report the predictions that failed with the same prominence as those that
held; three of the five produced at least one refuted sub-claim, and one of the
refutations (\cref{sec:app-chains}) overturns a claim made in our own earlier
internal digest.

%-----------------------------------------------------------------------------
\subsection{Serial investigation chains}
\label{sec:app-chains}

\paragraph{Prediction.} A multi-hop investigation --- alert, enrichment,
cross-reference, hypothesis, drill-down --- issues one request per hop, and each
hop blocks on the previous one. Concurrency is therefore~1 \emph{by
construction}: not a configuration we chose, but a property of the workload.
The mechanism of \cref{sec:mechanism} predicts that AR's batching multiplier is
structurally unavailable here, so the operative comparison is serial-versus-serial,
which the cost laws say diffusion wins at the natural output length of
${\sim}\naturalOutput$ tokens. It predicts further that per-hop latency is
constant in depth, so total latency is linear in depth, and --- because chain
latency is a \emph{sum} of per-hop latencies --- that diffusion's single-request
tail penalty should be attenuated by the central limit theorem as depth grows.

\paragraph{Measurement.} We run chains of depth $D \in \{1,3,5,10,20\}$, each
hop a single \texttt{generate} call that must complete before the next begins.

Per-hop latency is constant in depth for both models: the coefficient of
variation of the per-hop $p50$ across the five depths is \chainSlopeCVdg\% for
diffusion and \chainSlopeCVar\% for AR. Total chain latency is therefore linear
in $D$ for both, and the mean $p50$ ratio is stable at \chainMeanRatio$\times$
in diffusion's favour across every depth.

The tail behaves as predicted, and this is the result. At $D=1$ the
$p99.9$ ratio is \chainTailRatioDOne$\times$ --- diffusion's tail penalty almost
exactly cancels its mean advantage, so a single-hop, tail-sensitive deployment
sees no benefit at all. By $D=3$ the ratio has already migrated to
\needsnum{chainTailRatioDThree}$\times$, and it stays in that band
(\needsnum{chainTailRatioDFive}$\times$ at $D=5$,
\needsnum{chainTailRatioDTen}$\times$ at $D=10$,
\chainTailRatioDTwenty$\times$ at $D=20$), converging toward the
\chainMeanRatio$\times$ mean ratio rather than the single-request tail ratio.
Sequential composition averages the per-hop distributions, and the tail penalty
that dominates at $D=1$ washes out. The analytical CLT prediction agrees: with
per-hop $\sigma = \chainSigmaDG$~ms for diffusion and $\sigma = \chainSigmaAR$~ms
for AR, the predicted cumulative $p99.9$ at $D=20$ is
\needsnum{chainCumTailDGdTwenty}~ms versus
\needsnum{chainCumTailARdTwenty}~ms.

\paragraph{Correction: quality is not stable across depth.} An earlier internal
digest recorded this experiment's quality prediction as ``quality stable across
depths.'' Read against the raw per-depth artifact, that is false, and we state
the correction plainly. \emph{Both} models drop sharply between $D=1$ and $D=3$
--- diffusion from \needsnum{chainFOneDGdOne} to
\needsnum{chainFOneDGplateau}, AR from \needsnum{chainFOneARdOne} to
\needsnum{chainFOneARplateau} --- and only then plateau. The plateau is stable;
the transition into it is not, and both models settle well below their
single-hop level. Furthermore, AR is
approximately \needsnum{chainQualityGap}~F1 \emph{better} than diffusion at every
depth we measured, including $D=1$.

\paragraph{Verdict.} The chain result is a \emph{latency} result and nothing
more. The latency prediction --- linearity in depth, and CLT attenuation of the
tail penalty under composition --- is confirmed. The accompanying quality claim
is refuted: chained extraction degrades sharply for both architectures, and AR
holds a consistent quality lead throughout. Any deployment argument built on
this experiment must be a latency argument.

%-----------------------------------------------------------------------------
\subsection{Wide-schema normalization}
\label{sec:app-wide}

\paragraph{Prediction.} Widening the target schema widens the output. The cost
laws predict that AR's latency grows linearly in the number of fields $F$ while
diffusion's grows in canvas steps, so the latency ratio should grow with $F$.
A second, weaker prediction --- that joint denoising of all fields on one canvas
absorbs the $O(F^2)$ growth in pairwise consistency constraints --- gives a
falsifiable kill condition: diffusion's cross-field consistency must not fall
below $80\%$ at $F=50$.

\paragraph{Measurement.} We sweep $F \in \{12, 25, 50, 100\}$.

\begin{table}[t]
  \centering
  \small
  \caption{Wide-schema normalization. The latency ratio grows with schema
  width, but AR holds higher cross-field consistency at \emph{every} width.
  Campaign~B; \modelDG{} vs \modelARthree{}.}
  \label{tab:wide}
  \begin{tabular}{rccc}
    \toprule
    Fields $F$ & latency ratio AR/DG & diffusion consist.\ (\%) & AR consist.\ (\%)\\
    \midrule
    12  & \needsnum{wideRatioTwelve}     & \wideConsistTwelve  & \needsnum{wideConsistARTwelve}\\
    25  & \needsnum{wideRatioTwentyFive} & \needsnum{wideConsistTwentyFive} & \needsnum{wideConsistARTwentyFive}\\
    50  & \wideRatioFifty                & \wideConsistFifty   & \needsnum{wideConsistARFifty}\\
    100 & \needsnum{wideRatioHundred}    & \wideConsistHundred & \wideConsistARHundred\\
    \bottomrule
  \end{tabular}
\end{table}

The latency prediction holds: the ratio rises to \wideRatioFifty$\times$ at
$F=50$ and stays there. The kill condition is not triggered --- diffusion's
consistency plateaus at \wideConsistFifty--\wideConsistHundred\%, well above the
$80\%$ threshold, and the plateau is non-monotonic (consistency \emph{rises}
from \wideConsistTwelve\% to \needsnum{wideConsistTwentyFive}\% before falling).
Schema validity moves in the opposite direction from consistency, rising from
\needsnum{wideValidDGTwelve}\% at $F=12$ to \needsnum{wideValidDGHundred}\% at
$F=100$: the two metrics decouple, and validity is not a proxy for coherence.

\paragraph{Two caveats that bound the claim.} First, AR holds
\wideConsistARHundred\% consistency at $F=100$ against diffusion's
\wideConsistHundred\%, and AR is ahead at every width in \cref{tab:wide}.
\textbf{AR wins consistency at every schema width we measured.} The surviving
diffusion claim is that consistency degrades \emph{gracefully} and stays above
the kill threshold, not that it is superior. Second, diffusion emits more tokens
than AR at every width --- \wideTokensDGHundred{} versus \wideTokensARHundred{}
mean output tokens at $F=100$ --- so part of the measured latency ratio is
diffusion doing more work per record, and part of the consistency gap may itself
be a length effect rather than an architectural one.

\paragraph{Verdict.} Latency advantage confirmed and growing with $F$; graceful
degradation confirmed; superior consistency \emph{refuted}.

%-----------------------------------------------------------------------------
\subsection{Bounded-deadline completion}
\label{sec:app-deadline}

\paragraph{Prediction.} If diffusion's latency is content-independent and set by
the number of denoising steps, then a fixed step budget yields a complete
canvas --- a whole record or nothing --- whereas AR under a token budget derived
from its own cost law ($\texttt{max\_tokens} = D / \arSlope$) must stop wherever
the deadline falls, mid-record. The prediction is about \emph{completeness},
not quality: diffusion should emit syntactically complete records at deadlines
where AR emits truncated ones.

\paragraph{Measurement.} We sweep deadlines $D \in \{200, 600, 800, 1000, 1500,
2000\}$~ms. Diffusion (\modelLLaDA{}, $K=8$, single canvas, Campaign~D)
reaches \dmRecordComplete\% record completeness and \dmRecordComplete\% field
completeness at a $p50$ of \needsnum{dmDeadlinePFifty}~ms. AR's field
completeness rises with the deadline and saturates at \arFieldCompleteMax\%
--- eight of twelve fields --- while its record completeness is
\arRecordComplete\% at every deadline tested. AR's quality collapses at the
tightest deadlines as predicted, from \arDeadlineFOne{} F1 at $D \geq 1500$ down
to \needsnum{arDeadlineFOneTwoHundred}{} at $D=200$.

\paragraph{Open: this is not yet a clean architectural win.} AR's record
completeness is \arRecordComplete\% at \emph{every} deadline, including
$D=2000$, where it is given \needsnum{arDeadlineTokensTwoK} output tokens and
finishes in \needsnum{arDeadlinePFiftyTwoK}~ms --- comfortably inside its
budget. A model that finishes early with tokens to spare and still emits no
parseable complete record is more likely reporting a harness defect than an
architectural limit. We suspect a fenced-code preamble
(\verb|Output:```json|) consuming
part of the token budget before the record begins. This is filed as issue~\#38
and is unresolved. \open{\#38} Until it closes, the correct statement of this
result is: \emph{AR's record completeness was \arRecordComplete\% under our
generate-and-parse harness}, with the caveat attached. It is not evidence that
AR cannot complete a record under a deadline.

\paragraph{Quality runs the other way.} On the same experiment AR's Field-F1 is
\arDeadlineFOne{} against diffusion's \dmDeadlineFOne{}. AR produces incomplete
records of higher per-field quality; diffusion produces complete records of
lower quality. Whether that trade is worth taking is a property of the consumer
--- a downstream schema validator cares about completeness, a downstream
analyst cares about F1 --- and not a property of the architectures.

\paragraph{Verdict.} Completeness prediction passes under our harness but is
compromised by an open bug; quality prediction refuted.

%-----------------------------------------------------------------------------
\subsection{Iterative refinement}
\label{sec:app-refine}

\paragraph{Prediction.} Masked diffusion can re-mask an arbitrary subset of an
already-generated record and re-denoise only those positions, conditioned on the
positions it keeps. AR has no equivalent: correcting one field means
regenerating the record, or prompting the model to copy the good fields back
verbatim. The prediction is that diffusion improves across refinement rounds
driven by a critic that flags suspect fields, and that AR does not.

\paragraph{Measurement.} Three rounds; in each round after the first, a critic
selects fields to re-mask. Diffusion improves monotonically over the three
rounds for a total gain of \refineGainDM{} F1 (from
\needsnum{refineDMFirst} through \needsnum{refineDMSecond} to
\refineDMFinal). AR is flat: \refineGainAR{} F1 across the same three rounds,
which is within run-to-run noise of no change at all.

\paragraph{The bound on this claim.} Diffusion after three rounds
(\refineDMFinal) still trails AR after \emph{one} (\refineARFirst). The
refinement loop costs three model calls and closes less than half of the initial
gap. The defensible claim is therefore \emph{``diffusion has a refinement
mechanism that AR lacks,''} which the data supports, and not \emph{``diffusion
ends up better,''} which it does not. This experiment's paired negative result
--- the same masking mechanism failing outright at one-shot imputation --- is in
\cref{sec:neg-critic}, and the contrast between them is more informative than
either result alone.

\paragraph{Verdict.} Mechanism confirmed; outcome does not overtake AR.

%-----------------------------------------------------------------------------
\subsection{Priority preemption}
\label{sec:app-vip}

\paragraph{Prediction.} A latency-critical record --- a FATAL line arriving
mid-stream --- must be served inside a $1$~s SLA while background work
continues. The cost law says diffusion's service time is
$\dgBase$~ms and is independent of the record's content, so if the serving
engine can preempt, that service time is available on demand.

\paragraph{Measurement.} Both vLLM~0.25.1 preemption paths were verified live:
per-request \texttt{priority} together with
\texttt{--scheduling-policy priority}, with the engine observed preempting
lower-priority running sequences. Diffusion survives \texttt{priority} combined
with \texttt{--max-num-seqs 1} without triggering the mixed-phase crash
discussed in \cref{sec:two-walls}. Under background load the preempted
high-priority request completes at a lag $p50$ of \needsnum{vipLagPFiftyDG}~ms
and $p99$ of \needsnum{vipLagPNineNineDG}~ms.

\paragraph{Verdict.} Confirmed, with the scope it deserves: this is an
engineering check that the content-independent service time is reachable
mid-stream through a standard scheduler, not an independent measurement of an
architectural property.

%-----------------------------------------------------------------------------
\subsection{Scoreboard}
\label{sec:app-scoreboard}

\begin{table}[t]
  \centering
  \small
  \caption{Predictions derived from the cost laws, and what the measurements
  returned. Refutations and confirmations are listed together and in the order
  the sub-claims were made.}
  \label{tab:scoreboard}
  \begin{tabular}{p{0.44\linewidth}lp{0.30\linewidth}}
    \toprule
    Prediction & Verdict & The one number\\
    \midrule
    Per-hop chain latency constant in depth
      & Confirmed & slope CV \chainSlopeCVdg\% (DG)\\
    Tail penalty washes out under sequential composition
      & Confirmed & $p99.9$ ratio \chainTailRatioDOne$\times\!\to\!$\chainTailRatioDTwenty$\times$\\
    Chain quality stable across depth
      & \textbf{Refuted} & AR ahead by \needsnum{chainQualityGap} F1 at every depth\\
    Latency ratio grows with schema width
      & Confirmed & \wideRatioFifty$\times$ at $F=50$\\
    Consistency kill condition ($<80\%$ at $F=50$)
      & Not triggered & \wideConsistFifty\%\\
    Diffusion wins cross-field consistency
      & \textbf{Refuted} & AR \wideConsistARHundred\% at $F=100$\\
    Diffusion completes records under a deadline
      & Passes, \open{\#38} & \dmRecordComplete\% vs \arRecordComplete\%\\
    Diffusion wins quality under a deadline
      & \textbf{Refuted} & AR \arDeadlineFOne{} vs DM \dmDeadlineFOne\\
    Diffusion improves under critic-driven refinement
      & Confirmed & \refineGainDM{} F1 over three rounds\\
    Refinement lifts diffusion past AR
      & \textbf{Refuted} & \refineDMFinal{} vs \refineARFirst{} after one round\\
    Content-independent service time is preemptible
      & Confirmed & $\dgBase$~ms base service time\\
    \bottomrule
  \end{tabular}
\end{table}

\Cref{tab:scoreboard} collects the eleven sub-claims. Six are confirmed, one is
a kill condition that failed to fire, one passes under an open harness bug, and
four are refuted. The pattern in the refutations is consistent and worth naming:
every refuted sub-claim is a \emph{quality} claim, and every confirmed sub-claim
is a \emph{cost} claim. The cost laws predict cost well. They do not license
quality claims, and \cref{sec:negative} shows what happens when we tried to
extend them into quality territory anyway.
